% !TEX root = ../main_hessenberg.tex
\section{Formalization as the primary artifact}\label{sec:methodology}

\sloppy

In dependent type theory, a Lean formalization is inherently a proof~\cite{deMoura2021Lean4}, requiring no accompanying prose. On this basis, we treat the Lean development as the primary mathematical artifact of this project, and the prose as merely one rendering of it. Because the human contribution shifts from writing code to exercising mathematical judgment and steering the AI, the central object of study of this paper is precisely prompt design. Getting the proof to compile is the starting point of that work, not its endpoint. What distinguishes a primary mathematical artifact from a mere verification record lies entirely past that threshold: it depends on whether the code, in its own form, legibly conveys its mathematical content. Finally, we acknowledge that neither author has a formal software engineering background. We treat AI tools as informal mentors in software practice, working to elevate this project to structural standards we did not initially command, and we expect our execution will still reflect this learning curve.

\medskip
\noindent\textbf{Prompt example --- project bootstrap.}
\begin{quote}\itshape\small
The goal is not just proving the main theorem; it is producing code
whose main theorem signature reads as a literal translation of the paper
sentence. Engineering tax should be separated from mathematical content
and encapsulated in independent modules; every name should carry clear
mathematical semantics. Survey the paper's core mathematical objects;
propose a Foundation-structure design. Output the design only; do
not implement until I confirm.
\end{quote}

%\subsection{Engineering tax: see it, then move it}\label{ssec:tax}

We refer to type-theoretic overhead with no paper analogue as
\emph{engineering tax}. Left unmanaged, this tax sits mixed among
mathematical-layer declarations and muddles the mathematical narrative.
Three forms recur in this project:
\begin{itemize}[leftmargin=2em, itemsep=2pt, topsep=2pt]
\item bound-carrying boilerplate, where anonymous-constructor
  expressions bundle a value together with proofs that it lies in range;
\item index translations, where the formal system's natural index
  types differ from the paper's mathematical indexing conventions;
\item case analyses on small numeric ranges, where boundary cases
  dismissed by the paper in a single sentence require explicit,
  separate handling.
\end{itemize}
A typical instance is the lemma \texttt{cyclic\_shift\_arc\_forward}:
its statement reads like a standard mathematical theorem, while its proof body
spans roughly $80$ lines combining all three forms of overhead. A casual
reading of the code easily overlooks the true density of this burden.

To surface this hidden overhead, we tag every declaration as \texttt{Math}, 
\texttt{Eng}, or \texttt{Mixed}. Because these tags are formatted as docstrings 
above each declaration, they travel inherently with the code:
\begin{lstlisting}
/-! ## Mathematical layer -- paper-side vertex type. A `Vertex n`
    is a label `v ∈ {1, ..., n}` (Sec. 0 of the paper). -/

/-! ## Engineering layer -- Type B bridge: `Vertex n ↔ Fin n`
    (paper 1-indexed ↔ Lean 0-indexed). -/
\end{lstlisting}
We enforce this tagging discipline with a small project-specific linter, 
alongside Mathlib's standard tooling~\cite{The_mathlib_Community_2020} 
(naming, docstrings, style, plus \texttt{shake} for unused imports), 
ensuring that violations surface mechanically.

Engineering tax cannot be eliminated, but its location can be chosen.
We encapsulate it in independent modules so the mathematical layer
interacts with them strictly through clean APIs. This applies Parnas's
information-hiding principle~\cite{Parnas1972} to formalization: each
module hides one engineering concern from the mathematical layer. The
core of this effort is a small Foundation layer of bundled structures,
each encapsulating a paper-side mathematical object as a Lean
type~\cite{Baanen_2022, commelin2023abstractionboundariesspecdriven, pittphilsci13482}.
The layer introduces \texttt{Vertex}, \texttt{ActiveSet},
\texttt{Digraph}, \texttt{OrthogonalHessenberg},
\texttt{UnreducedAngles}, and others---allowing each structure to absorb the
corresponding combination of hypotheses just once. Three further modules
(\texttt{Singleton.Fin}, \texttt{EntryAccess},
\texttt{PartialProduct}) are factored out by thematic content.

Such structural refactors involve large-scale cascades---the largest spanned
$12$ files and roughly $232$ call sites---where the design judgment is
made by humans and the execution is carried out by AI. To illustrate the impact, consider the Bridge theorem signature before the 
structural pass:
\begin{lstlisting}
theorem digraph_model (hn : 2 ≤ n) (θ : Fin (n - 1) → ℝ)
    (hunred : ∀ k, Real.sin (θ k) ≠ 0) :
    ∀ i j : Fin n,
      givensProduct θ i j ≠ 0 ↔ Arc n (activeSet θ) (i.val + 1) (j.val + 1)
\end{lstlisting}
After the structural pass:
\begin{lstlisting}
theorem digraph_model (hn : 2 ≤ n) (θ : UnreducedAngles n) :
    θ.product.supportDigraph = θ.activeSet.digraph
\end{lstlisting}
The engineering tax has not disappeared; it simply now lives within the Foundation
modules. Consequently, the formal Bridge theorem reads exactly as the 
natural-language paper sentence.

\medskip
\noindent\textbf{Prompt example --- abstraction design.}
\begin{quote}\itshape\small
Scan all uses of \texttt{Fin n} that represent paper-side vertices.
Report the count of anonymous-mk constructions, the count of
$0$-indexed-to-$1$-indexed translations, and the most common usage
patterns. From the data, propose a \texttt{Vertex n} structure with
constructors, comparison, neighbouring, and a round-trip API to
\texttt{Fin n}. Output the design only.
\end{quote}

%\subsection{Letting mathematics speak}\label{ssec:speak}

Once the engineering tax is isolated, the mathematical layer can carry
semantics directly: names become clear mathematical statements, symbols
match the paper, and the main theorem signature reads as a literal
translation of its natural-language counterpart.

For naming, we apply a simple heuristic: read the name aloud and check
whether it forms a complete mathematical thought. A name that fails
this test usually signals that the declaration's role has not been fully
resolved---it may be too specific, too abstract, or in the wrong place. 
For instance, \texttt{extended\_diag\_isPlusMinusOne} fails because it 
leaves the subject ambiguous; revised to 
\texttt{boundaryAppended\_diag\_isPlusMinusOne}, it forms a clear sentence. 
This test serves both human and AI readability. Placeholder names like 
\texttt{aux\_lemma\_3} leave the AI with no semantic basis for judgment in 
search and completion, whereas a name that encodes its own statement exposes 
the underlying mathematics directly to the model.

We use the form of the main theorem itself as the project's completion 
criterion: when the apex signature reads exactly as the paper sentence, 
the work is done. A green build is necessary, but not sufficient. Consider the apex theorem immediately after proof completion, before the 
engineering tax has been extracted:
\begin{lstlisting}
theorem classification {n : ℕ} (hn : 2 ≤ n) :
    (∀ S : Finset ℕ, ValidS n S → ∃ θ : Fin (n-1) → ℝ,
        (∀ k, Real.sin (θ k) ≠ 0) ∧ activeSet θ = S) ∧
    (∀ θ : Fin (n-1) → ℝ, (∀ k, Real.sin (θ k) ≠ 0) →
        ∀ i j : Fin n, givensProduct θ i j ≠ 0 ↔
          Arc n (activeSet θ) (i.val + 1) (j.val + 1)) ∧
    (∀ S S' : Finset ℕ, ValidS n S → ValidS n S' →
        2 ≤ S.card → 2 ≤ S'.card → DigraphIso n S S' → S = S')
\end{lstlisting}
Here, engineering overhead sits alongside mathematical content in the
signature, forcing the reader to mentally filter it out. The code is 
mathematically verified, but it has not assumed responsibility for legibility. Contrast this with the signature after the structural pass:
\begin{lstlisting}
theorem classification {n : ℕ} (hn : 2 ≤ n) :
    (∀ S : ActiveSet n, ∃ θ : UnreducedAngles n, θ.activeSet = S) ∧
    (∀ θ : UnreducedAngles n, θ.product.supportDigraph = θ.activeSet.digraph) ∧
    (∀ S S' : ActiveSet n, 2 ≤ S.elements.card → 2 ≤ S'.elements.card →
        S.digraph ≅ S'.digraph → S = S') :=
  ⟨completeness, digraph_model hn, fun S S' => rigid S S' hn⟩
\end{lstlisting}
Every conceptual object in the signature is now packaged inside its own
structure. The entire declaration reads as a direct translation of
Theorem~\ref{thm:classification}: for $n \ge 2$, every active set
arises from some unreduced angle vector; the support digraph of the
Givens product equals the one given by the active set; and for
$|S|, |S'| \ge 2$, digraph isomorphism implies $S = S'$.

Both signatures encode identical mathematical content; the difference
lies entirely in whether the code assumes responsibility for its own legibility. 
Only when the formal development is treated as the primary mathematical 
artifact can this standard be posed and met.

\medskip
\noindent\textbf{Prompt example---naming review.}
\begin{quote}\itshape\small
Apply the natural-language reading test to every helper in this file.
For each name that fails, identify the failure mode (missing subject,
placeholder word, unexpanded abbreviation) and propose a revised name
that states what the object is, what it is built from, and what
property it has. Report the rename list; confirm before implementing.
\end{quote}
